% ============================================================================
% ABSTRACT
% ============================================================================

\vspace*{2cm}

\begin{center}
\textbf{Shape-Height Transform for DNA: \\[0.3cm] Mathematical Framework and Comprehensive Validation of Three Orthogonal Geometric Units}
\end{center}

\vspace{1cm}

DNA structure fundamentally determines biological function, yet existing tools for DNA characterization focus primarily on helical parameters (twist, roll, tilt) rather than cross-sectional geometry. We present the Shape-Height Transform (SHT), a novel mathematical framework for quantifying local DNA geometry at basepair resolution through three orthogonal geometric units: Vij (curvature), Shambhian (cross-sectional flatness), and Mamton (surface roughness).

\vspace{0.5cm}

We validate each unit rigorously on 202 independent PDB structures, demonstrating:
\begin{enumerate}
    \item \textbf{Mathematical correctness}: Each unit accurately measures its geometric property with quantitative agreement to known structures (nucleosome radius, A-tract rigidity, protein deformation).
    
    \item \textbf{New discovery}: Vij and Mamton exhibit strong coupling ($r = 0.83$, $p < 10^{-16}$), revealing a quantitative bending-roughness relationship supporting the hypothesis that differential strain from bending induces cross-sectional asymmetry.
    
    \item \textbf{Functional prediction}: A Random Forest classifier trained on the three geometric units achieves 91\% accuracy in predicting DNA functional state (nucleosome, protein-bound, free duplex, non-canonical), with blind validation on the nucleosome structure 1F66 confirmed by independent PDB metadata.
\end{enumerate}

\vspace{0.5cm}

This monograph provides complete mathematical derivations, computational implementation details, comprehensive validation across functional classes, and biological insights. All claims are backed by quantitative evidence ($p < 0.01$), with effect sizes exceeding 0.8 (Cohen's $d$). We establish SHT as a rigorous framework enabling geometry-based functional prediction from DNA structures alone, with applications to nucleosome positioning, drug design, and DNA nanotechnology.

\vspace{0.5cm}

\noindent\textbf{Keywords:} DNA geometry, shape descriptors, cross-sectional analysis, curvature, functional classification, structural bioinformatics

\vspace{2cm}
